\documentclass[11pt,a4paper]{article}
\usepackage[T1]{fontenc}
\usepackage[utf8]{inputenc}
\usepackage{lmodern}
\usepackage{amsmath,amssymb,amsthm,mathtools,mathrsfs}
\usepackage{graphicx,float,xcolor,multicol}
\usepackage[shortlabels]{enumitem}
\usepackage[margin=27mm]{geometry}
\usepackage{microtype,needspace,placeins}
\usepackage{tikz,tikz-cd}
\usetikzlibrary{arrows.meta,calc,positioning,patterns,decorations.pathreplacing}
\usepackage[colorlinks=true,linkcolor=teal,urlcolor=teal,citecolor=teal]{hyperref}
\setlength{\parindent}{0pt}
\setlength{\parskip}{.6em}
\setcounter{secnumdepth}{0}
\allowdisplaybreaks
\raggedbottom
\providecommand{\cross}{\times}
\DeclareUnicodeCharacter{2020}{\ensuremath{\dagger}}
\DeclareMathOperator{\supp}{supp}
\DeclareMathOperator*{\esssup}{ess\,sup}
\providecommand{\chapter}{\section}
\newenvironment{tool}[2]{\subsection*{Tool #1 --- #2}}{}
\newcommand{\ToolRef}[1]{Tool #1}
\newcommand{\FigureTag}[2]{\textit{Figure #1}}
\newcommand{\Result}[1]{\subsection*{#1}}
\newcommand{\Application}[1]{\subsubsection*{Application\if\relax\detokenize{#1}\relax\else\space--- #1\fi}}
\newcommand{\ProofHeading}{\par\textit{Proof.}\space}
\newcommand{\ProofOf}[1]{\par\textit{Proof of #1.}\space}
\newcommand{\RemarkHeading}{\par\textbf{Remark.}\space}
\newcommand{\NoteHeading}{\par\textbf{Note.}\space}
\newcommand{\ExerciseHeading}{\par\textbf{Exercise.}\space}

\graphicspath{{figures/surgeries/}}
\begin{document}
\section*{The Multiplier Map}
% BLOG-CONTENT-BEGIN

    
    
    The argument below is developed in greater detail in \cite{branner2014quasiconformal}, \cite{MR812271}, and \cite{MR762431}. We shall prove the following theorem. 
    
    \subsection*{Theorem 5.1 (Multiplier Map for a Unicritical Polynomial)}
    Let $\Omega$ be a \BlogPost{the-quadratic-family}{hyperbolic component} of $\mathcal{M}_d$ for $d\geq2$, and let $c^*$ be its centre. Then the multiplier map
    $$\Lambda_d:\Omega\longrightarrow\mathbb{D}$$
    is a proper holomorphic map of degree $d-1$, with $\Lambda_d^{-1}(0)=\{c^*\}$ and local degree $d-1$ at $c^*$. Equivalently,
    $$\Lambda_d:\Omega\backslash\{c^*\}\longrightarrow\mathbb{D}^*$$
    is an unbranched $(d-1)$-sheeted covering.

    Here the unicritical family of degree $d$ is parametrized by polynomials of the form $Q_c(z)=z^d+c$.
    
    
    
    Let $\mathbb{D}^*$ be the punctured disk, and let $\Omega^*=\Lambda_d^{-1}(\mathbb{D}^*)=\Omega\backslash\{c^*\}$. Sullivan proved that $\Lambda_d:\Omega^*\rightarrow\mathbb{D}^*$ is a covering map. Gleason gave an argument that helped in calculating the degree of that map.
    
    

    Now it remains to prove that the degree of the multiplier map $\Lambda_d:\Omega\rightarrow\mathbb{D}$ is $d-1$. If $\Omega$ has period $p$, then $Q_{c^*}^{\circ p}(0)=0$. The following lemma of Gleason completes the proof.
    
    \subsection*{Lemma 5.2}
    Consider the polynomials $G_n(c)=Q_c^{\circ n}(0)$ for $n\geq1$. All roots of $G_n$ are simple.
    
    
    \textit{Proof.}: We have,
    
    $$G_1(c) = c$$
    $$G_n(c) = (G_{n-1}(c))^d + c$$
    
    It follows that $G_n$ is monic with integer coefficients. The roots of $G_n$ are therefore algebraic integers over $\mathbb{Z}$. Assume that there exists an $n>1$ and a root $\tilde{c}$ of $G_n$ which is not simple. Then $\tilde{c}$ must be a root of both $G_n$ and its derivative. That is,
    $$G_n'(c) = d(G_{n-1}(c))^{d-1}G_{n-1}'(c)+1$$
    
    It follows that,
    $$(G_{n-1}(\tilde{c}))^{d-1}G_{n-1}'(\tilde{c}) = \frac{-1}{d}$$
    
    Since the set of algebraic integers $\mathbb{A}$ forms a ring, we have from the above equation that $\frac{-1}{d}$ is an algebraic integer. But since $\mathbb{Q}\cap\mathbb{A} = \mathbb{Z}$, $G_n$ has no multiple roots.


    Note that for any $c\in \Omega$ the points in the attracting $p$-cycle are solutions of
    $$F(c,z) = Q_c^{\circ p}(z) - z = 0.$$

    The map $F$ satisfies the hypothesis of the Implicit Function Theorem at the point $(c^*,0)$, since
    $$F_z(c^*,0)=(Q_{c^*}^{\circ p})'(0)-1=-1,
    \qquad F_c(c^*,0)=G_p'(c^*)\neq0$$
    by Gleason's lemma. So, in a neighbourhood of $c^*$ the points in the attracting $p$-cycle can therefore be expressed as functions of $c$. Let $\alpha_0(c)$ denote the periodic point in the Fatou component containing the critical point $0$, so $\alpha_0(c^*)=0$, and
    $$\alpha_0'(c^*)=-\frac{F_c(c^*,0)}{F_z(c^*,0)}=G_p'(c^*)\neq0,$$
    so $\alpha_0(c)$ has a simple zero at $c=c^*$. Set $\alpha_j(c)=Q_c^{\circ j}(\alpha_0(c))$ for $j=1,\ldots,p-1$. The multiplier map can be written as
    $$\Lambda_d(c)=(Q_c^{\circ p})'(\alpha_0(c))
    =\prod_{j=0}^{p-1}Q_c'(\alpha_j(c))
    =d^p\prod_{j=0}^{p-1}\alpha_j(c)^{d-1}.$$

    Since $\alpha_0$ has a simple zero at $c^*$ and $\alpha_j(c^*)\neq0$ for $1\leq j<p$,
    $$\operatorname{ord}_{c^*}\Lambda_d
    =(d-1)\sum_{j=0}^{p-1}\operatorname{ord}_{c^*}\alpha_j
    =d-1.$$
    Thus the local degree of $\Lambda_d$ at $c^*$ is $d-1$. Since $\Lambda_d:\Omega^*\rightarrow\mathbb{D}^*$ is a covering, it has $d-1$ sheets. With this we conclude the proof of Theorem $5.1$.
% BLOG-CONTENT-END
\bibliographystyle{amsalpha}
\bibliography{tools/profiles/surgeries/MSc_Dissertation}
\end{document}
